% ============================================================
% Chapter 5: Development
%   Section: News Acquisition: Discovery and Extraction
%     Subsection: The Source Catalogue
% Included from chapters/05-development/02-news-acquisition/news-acquisition.tex
% ============================================================

\subsection{The Source Catalogue}
\label{subsec:scraping-sources}

Each news source is one YAML file in \path{backend/data/sources/},
loaded by \texttt{SourceRepository} into a \texttt{NewsSource} model.
Only an identifier, a name and a base URL are required; the fields
that matter to acquisition are listed in
Table~\ref{tab:source-fields}.

\begin{table}[ht]
\centering
\small
\caption{Fields of a source that acquisition reads.}
\label{tab:source-fields}
\begin{tabularx}{\linewidth}{@{}lX@{}}
\toprule
Field & Use \\
\midrule
\texttt{rss\_url}             & The feed read first by discovery. \\
\texttt{language}             & \texttt{en} or \texttt{es}; selects the
                                 keyword lists and section names used to
                                 filter its links. \\
\texttt{groups}               & The topic groups the source has real
                                 sections for
                                 (Section~\ref{subsec:scraping-groups}). \\
\texttt{reliability\_index}   & Editorial reliability in $[0,1]$, used
                                 when its pages are ranked as evidence
                                 and when its candidates are ranked. \\
\texttt{enabled}              & Whether the source is ingested at all. \\
\texttt{requires\_javascript} & Sends its articles to the browser first
                                 (Section~\ref{subsec:scraping-cascade}). \\
\texttt{positive\_editorial}  & The outlet publishes only positive news;
                                 the mission screen does not apply to it. \\
\texttt{metadata.selectors}   & Optional CSS selectors for the title,
                                 body, author and date, for layouts no
                                 generic parser understands. \\
\bottomrule
\end{tabularx}
\end{table}

\paragraph{Composition.} The catalogue holds 36 sources, 20 in English
and 16 in Spanish: 26 news outlets, four public bodies (NASA, the
European Space Agency, the World Health Organization and UN News), four
academic publishers (MIT News, \textit{Nature}, The Conversation and
Yale Environment 360) and two blogs. Three of them (Good News Network,
Positive News and Reasons to be Cheerful) publish only positive news.
Appendix~\ref{app:sources} lists them all.

\paragraph{Enabled and disabled.} 32 sources are enabled, 19 in English
and 13 in Spanish. A disabled source is not ingested, probed or
recognised from a posted URL, but it is still read by the reliability
map, so its pages keep their rating when they are found as evidence.
That is the mechanism for rating a domain without ingesting it, and it
is used for two different reasons. The fact-checkers Newtral and
Maldita are disabled because a fact-check quotes the claim it
debunks, and claim extraction would take that quotation for the
article's own claim. Agencia EFE and Reuters are disabled because they
refused every request from 23~September to 6~October~2026 (HTTP 403
and 401 respectively).

\paragraph{One source per domain.} Reliability and the recognition of
posted URLs are both looked up by domain, so a second file on the same
domain (BBC Mundo on \texttt{bbc.com}, for instance) would silently
override the first. A test of the repository fails on any such pair.

\paragraph{How a source is admitted.} The catalogue grew from 12 to 36
sources on 28~September~2026, and every addition passed the same test
before its file was committed: its feed had to produce article links,
and two of those articles had to extract through the real cascade.
Seven candidates failed it: NPR, Phys.org and El Confidencial (the
feed worked, the articles would not extract), ScienceDaily, DW Español
and CSIC (no article links in the feed) and the Associated Press
(HTTP 403). RTVE's dead feed was replaced the same day.
